\documentclass{report}
\usepackage{amsmath,amssymb}
\setlength{\parindent}{0mm}
\setlength{\parskip}{1em}
\begin{document}
\begin{center}
\rule{6in}{1pt} \
{\large Benjamin, J Sturdevant \\
{\bf Iterative Solution Method for an Implicit Orbit Averaged Particle-in-Cell Model}}

8901 Grant St Apt 336 \\ Thornton \\ CO 80229
\\
{\tt benjamin.sturdevant@colorado.edu}\\
Scott Parker\end{center}

Present kinetic simulations of turbulence in magnetized plasmas employ
models from gyrokinetic theory, which is based on a number of ordering
assumptions used to reduce the Vlasov-Maxwell system to eliminate high
frequency phenomena. Recently, a second order accurate, implicit
particle-in-cell (PIC) model using ions which evolve according to the
exact Newton-Lorentz equations of motion has been developed to study
low-frequency phenomena in magnetized plasmas. This model promises to be
physically reliable in regimes where the ordering assumptions of
gyrokinetics may be in question, but presents some additional
computational challenges to produce numerically stable results while
retaining important kinetic effects.

This talk will focus on orbit averaging algorithms applied to the Lorentz
ion model. Orbit averaging is a time stepping method based on advancing
particles over several micro time steps for each macro time step over
which the fields are advanced. An implicit scheme is applied to the orbit
averaging algorithm to eliminate a time stepping constraint due to a high
frequency compressional Alfven wave present in the Lorentz ion model. The
implicit algorithm requires the solution of a large system of nonlinear
equations at each macro-time step. A solution method for this system
based on Picard iterations will be presented. An "inner" iteration is
used for the particle advances at the micro time steps and an "outer"
iteration in which the particle system and the field equations are
alternately updated is used to provide a self-consistent coupling at the
macro time step. A preconditioner may provide a means of increasing the
rate of convergence in this iterative process. Numerical results will be
presented along with a discussion on the convergence of the iterative
method.

[1] Y. Chen, S.E. Parker, Phys. Plasmas 16 (2009)
[2] J. Cheng, S.E. Parker, Y. Chen, D. Uzdensky, J. Comput. Phys. 245 (2013), 364


\end{document}
